\subsubsection*{1.4.6.}
后面将考虑如下情形：$X_\eta$ 局部紧，$X_0$ 由 $X_\eta$ 将一个紧子集 $V\subset X_\eta$ 压缩为点 $0$ 而得，$\mathrm{sp}$ 为压缩映射，且 $\pi_1(D^*)$ 的作用在 $V$ 外平凡。对 $\mathrm{sp}$ 的固有基变换于是给出，对 $K\in\operatorname{Ob}D^+(\mathcal X^*)$，
\begin{equation}\tag{1.4.6.1}
R\Psi_\eta(\beta^*K)_0=R\Gamma(V,K)
\end{equation}
以及
\begin{equation}\tag{1.4.6.2}
R\Gamma_{\{0\}}R\Psi_\eta(\beta^*K)
=R\Gamma_V(X_\eta,K),
\end{equation}
后者为带支集的上同调。

对 $\mathcal X$ 上的 $K$，假设
\[
\mathrm{sp}:\mathrm{sp}^*K_0\longrightarrow K_\eta
\]
在 $V$ 之外是拟同构。于是 $R\Phi(K)=0$ 在 $0$ 外成立。特别地，对 $\mathcal X$ 上的常值层 $\Lambda$，
\[
R^i\Phi(\Lambda)=\widetilde H^i(V,\Lambda),
\qquad\text{支集位于 }0,
\]
其中右端为约化上同调。

\subsubsection*{1.4.7.}
我们关心的是 $X_\eta$ 为可微流形，而 $V\subset X_\eta$ 为具有边界及内部 $V^\circ$ 的子流形的情形。此时 $X-\{0\}$ 是将 $\beta^*X_\eta$ 与 $D$ 上以 $X_\eta-V$ 为纤维的常值丛粘合得到的可微流形。映射
\begin{equation}\tag{1.4.7.1}
H_c^i(V^\circ,\Lambda)
\xrightarrow{\ \sim\ }H_V^i(X_\eta,\Lambda)
\end{equation}
是同构。这些群也记为 $H^i(V\bmod\partial V,\Lambda)$。同构 (1.4.7.1) 定义一个同构
\begin{equation}\tag{1.4.7.2}
R\Gamma_c(V^\circ,\Lambda)
\xrightarrow{\ \sim\ }
R\Gamma_{\{0\}}R\Psi_\eta(\Lambda),
\end{equation}
按照 XIII, 2.4.5 的步骤并通过 (1.4.7.2)，\underline{变差} 定义为一个以 $V$ 的约化上同调为源的态射：
\begin{equation}\tag{1.4.7.3}
\operatorname{Var}:
R\widetilde\Gamma(V,\Lambda)
\longrightarrow R\Gamma_c(V^\circ,\Lambda).
\end{equation}

$T$ 在对偶 $(V,\partial V)$ 上诱导的自同构仍记为 $T$，其在 $\partial V$ 上的限制 $\partial T$ 是恒等映射。变差 (1.4.7.3) 应当等于附属于此自同构 $T$ 的经典变差态射（3.1.9）。3.1.10 将证明这种形式的断言。

